\Needspace{6\baselineskip}\Article{347}
\textsc{Satz.} \textit{Ist \(F=\varphi(t,\allowbreak u.v\ldots)\) eine symmetrische\HistoricalNote{*)}{\normalfont Bekanntlich heißen diejenigen Funktionen symmetrisch, in denen alle Unbestimmten auf dieselbe Weise vorkommen, oder deutlicher, die sich nicht ändern, wie man auch die Unbestimmten untereinander vertauscht; solche Funktionen sind \textit{zum Beispiel} die Summe aller Unbestimmten, ihr Produkt, die Summe ihrer paarweisen Produkte usw.} algebraische} \textit{ganzrationale Funktion von \(f\) Unbestimmten \(t,\allowbreak u,\allowbreak v\) usw., und führt man nach Einsetzen der \(f\) in der Periode \((f,\allowbreak \lambda)\) enthaltenen Wurzeln für diese Unbestimmten den Wert von \(F\) nach den Regeln des Artikels 340 auf die Form} \[A+A^{\prime}[1]+A^{\prime\prime}[2]+\text{usw.}=W\] \textit{zurück, so haben die Wurzeln, die in diesem Ausdruck zu einer und derselben beliebigen Periode mit \(f\) Gliedern gehören, gleiche Koeffizienten.}

\textit{Beweis.} Es seien \([p],\allowbreak [q]\) zwei Wurzeln, die zu einer und derselben Periode gehören, und \(p,\allowbreak q\) seien positiv und kleiner als \(n\). Dann ist zu zeigen, dass \([p]\) und \([q]\) in \(W\) denselben Koeffizienten haben. Es sei \(q\equiv pg^{\nu e}\ (\mathrm{mod.}\ n)\); ferner seien die in \((f,\allowbreak \lambda)\) enthaltenen Wurzeln \([\lambda],\allowbreak [\lambda^{\prime}],\allowbreak [\lambda^{\prime\prime}]\) usw., wobei wir die Zahlen \(\lambda,\allowbreak \lambda^{\prime},\allowbreak \lambda^{\prime\prime}\) usw. als positiv und kleiner als \(n\) voraussetzen. Schließlich seien die kleinsten positiven Reste der Zahlen \(\lambda g^{\nu e},\allowbreak \lambda^{\prime} g^{\nu e},\allowbreak \lambda^{\prime\prime} g^{\nu e}\) usw. nach dem Modul \(n\) die Zahlen \(\mu,\allowbreak \mu^{\prime},\allowbreak \mu^{\prime\prime}\) usw.; diese sind offenbar mit den Zahlen \(\lambda,\allowbreak \lambda^{\prime},\allowbreak \lambda^{\prime\prime}\) usw. identisch, wenn auch in anderer Reihenfolge. Aus Artikel 340 folgt nun, dass sich \[\varphi([\lambda g^{\nu e}],[\lambda^{\prime} g^{\nu e}],[\lambda^{\prime\prime} g^{\nu e}]\ldots)=(\mathrm I)\] auf \[\begin{gathered}
A+A^{\prime}[g^{\nu e}]+A^{\prime\prime}[2g^{\nu e}]+\mathrm{usw.}\\
\text{oder auf}\quad A+A^{\prime}[\theta]+A^{\prime\prime}[\theta^{\prime}]+\mathrm{usw.}=(W^{\prime})
\end{gathered}\] zurückführen lässt, wobei \(\theta,\allowbreak \theta^{\prime}\) usw. die kleinsten Reste der Zahlen \(g^{\nu e},\allowbreak 2g^{\nu e}\) usw. nach dem Modul \(n\) bezeichnen. Daraus ist klar, dass \([q]\) in \((W^{\prime})\) denselben Koeffizienten hat wie \([p]\) in \((W)\). Man sieht aber ohne Mühe, dass die Entwicklung des Ausdrucks \((\mathrm I)\) dasselbe ergibt wie die Entwicklung von \(\varphi([\mu],\allowbreak [\mu^{\prime}].[\mu^{\prime\prime}]\ \text{usw.})\), da \(\mu\equiv\lambda g^{\nu e},\allowbreak \mu^{\prime}\equiv\lambda^{\prime} g^{\nu e}\) usw. \(\ (\mathrm{mod.}\ n)\) gilt. Dieser Ausdruck ergibt seinerseits dasselbe wie \(\varphi([\lambda],\allowbreak [\lambda^{\prime}],\allowbreak [\lambda^{\prime\prime}]\ \text{usw.})\), weil die Zahlen \(\mu,\allowbreak \mu^{\prime},\allowbreak \mu^{\prime\prime}\) usw. sich von \(\lambda,\allowbreak \lambda^{\prime},\allowbreak \lambda^{\prime\prime}\) usw. nur in der Reihenfolge unterscheiden, auf die es bei einer symmetrischen Funktion nicht ankommt. Daraus folgt, dass \(W^{\prime}\) mit \(W\) vollständig identisch ist; daher hat die Wurzel \([q]\) in \(W\) denselben Koeffizienten wie \([p]\). \textit{Q. E. D.}

Daraus ist klar, dass sich \(W\) auf die Form
\SourceLeaf{432}
\[A+a(f,1)+a^{\prime}(f,g)+a^{\prime\prime}(f,gg)\ldots+a^\varepsilon(f,g^{e-1})\] zurückführen lässt, wobei die Koeffizienten \(A,\allowbreak a\ldots a^\varepsilon\) feste Größen sind, die außerdem ganze Zahlen sind, wenn alle rationalen Koeffizienten in \(F\) ganze Zahlen sind. — So führt sich \textit{zum Beispiel}, wenn \(n=19\), \(f=6,\allowbreak \ \lambda=1\) gilt und die Funktion \(\varphi\) die Summe der paarweisen Produkte der Unbestimmten bezeichnet, ihr Wert auf \(3+(6,1)+(6,4)\) zurück.

Ferner sieht man leicht: Setzt man danach für \(t,\allowbreak u,\allowbreak v\) usw. die Wurzeln aus einer anderen Periode \((f,k\lambda)\) ein, so wird der Wert von \(F\) zu \[A+a(f,k)+a^{\prime}(f,kg)+a^{\prime\prime}(f,kgg)+\text{usw.}\]

\Needspace{6\baselineskip}\Article{348}
In jeder Gleichung \[x^f-\alpha x^{f-1}+\beta x^{f-2}-\gamma x^{f-3}\ldots=0\] sind die Koeffizienten \(\alpha,\allowbreak \beta,\allowbreak \gamma\) usw. symmetrische Funktionen der Wurzeln, nämlich \(\alpha\) die Summe aller Wurzeln, \(\beta\) die Summe ihrer paarweisen Produkte, \(\gamma\) die Summe ihrer Produkte zu je dreien usw. Daher ist in der Gleichung, deren Wurzeln die in der Periode \((f,\lambda)\) enthaltenen Wurzeln sind, der erste Koeffizient \(=(f,\lambda)\), während sich jeder der übrigen auf eine Form wie \[A+a(f,1)+a^{\prime}(f,g)\ldots+a^\varepsilon(f,g^{e-1})\] bringen lässt, wobei alle \(A,\allowbreak a,\allowbreak a^{\prime}\) usw. ganze Zahlen sind. Außerdem ist klar, dass die Gleichung, deren Wurzeln die in irgendeiner anderen Periode \((f,k\lambda)\) enthaltenen Wurzeln sind, aus jener hervorgeht, wenn man in jedem Koeffizienten \((f,1)\) durch \((f,k)\) ersetzt; ebenso \((f,g)\) durch \((f,kg)\) und allgemein \((f,p)\) durch \((f,kp)\). Auf diese Weise kann man also \(e\) Gleichungen \(z=0,\allowbreak z^{\prime}=0,\allowbreak z^{\prime\prime}=0\) usw. angeben, deren Wurzeln die in \((f,1)\), in \((f,g)\), \((f,gg)\) usw. enthaltenen Wurzeln sind, sobald die \(e\) Summen \((f,1),\allowbreak (f,g),\allowbreak (f,gg)\) usw. bekannt sind oder vielmehr, sobald \textit{eine} beliebige von ihnen gefunden ist, da sich nach Artikel 346 aus einer alle übrigen rational ableiten lassen. Damit ist zugleich die Funktion \(X\) in \(e\) Faktoren vom Grad \(f\) zerlegt; denn das Produkt aus den Funktionen \(z,\allowbreak z^{\prime},\allowbreak z^{\prime\prime}\) usw. ist offenbar \(=X\).

\textit{Beispiel.} Für \(n=19\) ist die Summe aller Wurzeln in der Periode \((6,1)\) \(=(6,1)=\alpha\); die Summe ihrer paarweisen Produkte ist \(=3+(6,1)+(6,4)=\beta\); ebenso
\SourceLeaf{433}
findet man für die Summe der Produkte zu je dreien \(=2+2(6,1)+(6,2)=\gamma\), für die Summe der Produkte zu je vieren \(=3+(6,1)+(6,4)=\delta\), für die Summe der Produkte zu je fünfen \(=(6,1)=\varepsilon\) und für das Produkt aller Wurzeln \(=1\). Daher umfasst die Gleichung \[z=x^6-\alpha x^5+\beta x^4-\gamma x^3+\delta xx-\varepsilon x+1=0\] sämtliche in \((6,1)\) enthaltenen Wurzeln. Setzt man nun in den Koeffizienten \(\alpha,\allowbreak \beta,\allowbreak \gamma\) usw. für \((6,1),\allowbreak (6,2),\allowbreak (6,4)\) jeweils \((6,2),\allowbreak (6,4),\allowbreak (6,1)\) ein, so erhält man die Gleichung \(z^{\prime}=0\), welche die Wurzeln in \((6,2)\) umfasst; wendet man dieselbe Vertauschung hier nochmals an, so erhält man die Gleichung \(z^{\prime\prime}=0\), welche die Wurzeln in \((6,4)\) umfasst, und das Produkt \(zz^{\prime}z^{\prime\prime}\) ist \(=X\).

\Needspace{6\baselineskip}\Article{349}
Meist ist es bequemer, besonders wenn \(f\) eine große Zahl ist, die Koeffizienten \(\beta,\allowbreak \gamma\) usw. nach dem Newtonschen Satz aus den Potenzsummen der Wurzeln abzuleiten. Es ist nämlich unmittelbar klar, dass die Summe der Quadrate der in \((f,\lambda)\) enthaltenen Wurzeln \(=(f,2\lambda)\), die Summe ihrer Kuben \(=(f,3\lambda)\) usw. ist. Schreiben wir also der Kürze wegen für \((f,\lambda),\allowbreak (f,2\lambda),\allowbreak (f,3\lambda)\) usw. \(q,\allowbreak q^{\prime},\allowbreak q^{\prime\prime}\) usw., so gilt \[\alpha=q,\quad 2\beta=\alpha q-q^{\prime},\quad 3\gamma=\beta q-\alpha q^{\prime}+q^{\prime\prime}\quad\text{usw.}\] Dabei sind die Produkte aus zwei Perioden nach Artikel 345 sofort in Summen von Perioden umzuwandeln. So werden in unserem Beispiel, wenn wir für \((6,1),\allowbreak (6,2),\allowbreak (6,4)\) jeweils \(p,\allowbreak p^{\prime},\allowbreak p^{\prime\prime}\) schreiben, die Größen \(q,\allowbreak q^{\prime},\allowbreak q^{\prime\prime},\allowbreak q^{\prime\prime\prime},\allowbreak q^{\prime\prime\prime\prime},\allowbreak q^{\prime\prime\prime\prime\prime}\) jeweils \(=p,\allowbreak p^{\prime},\allowbreak p^{\prime},\allowbreak p^{\prime\prime},\allowbreak p^{\prime},\allowbreak p^{\prime\prime}\); daher \[\begin{aligned}
\alpha&=p,\quad2\beta=pp-p^{\prime}=6+2p+2p^{\prime\prime}\\
3\gamma&=(3+p+p^{\prime\prime})p-pp^{\prime}+p^{\prime}=6+6p+3p^{\prime}\\
4\delta&=(2+2p+p^{\prime})p-(3+p+p^{\prime\prime})p^{\prime}\\
&\qquad+pp^{\prime}-p^{\prime\prime}=12+4p+4p^{\prime\prime}\quad\mathrm{usw.}
\end{aligned}\] Es genügt übrigens, nur die Hälfte der Koeffizienten auf diese Weise zu berechnen. Denn es lässt sich ohne Schwierigkeit beweisen, dass die letzten in umgekehrter Reihenfolge entweder den ersten gleich sind, nämlich der letzte \(=1\), der vorletzte \(=\alpha\), der davor \(=\beta\) usw., oder aus diesen jeweils hervorgehen, wenn man für \((f,1),\allowbreak (f,g)\) usw. \((f,-1),\allowbreak (f,-g)\) usw. oder \((f,n-1)\), \((f,n-g)\) usw. einsetzt. Der erste Fall tritt ein, wenn \(f\) gerade ist, der zweite, wenn \(f\) ungerade ist; der letzte Koeffizient ist aber stets \(=1\). Dies beruht auf dem Satz in Artikel 79; doch halten wir uns der Kürze wegen bei diesem Gegenstand nicht auf.
\SourceLeaf{434}

\clearpage
\Needspace{6\baselineskip}\Article{350}
\textsc{Satz.} \textit{Es sei \(n-1\) das Produkt aus drei positiven ganzen Zahlen \(\alpha,\allowbreak \beta,\allowbreak \gamma\); die} \textit{Periode \((\beta\gamma,\lambda)\), die \(\beta\gamma\) Glieder hat, bestehe aus \(\beta\) kleineren Perioden mit \(\gamma\) Gliedern,} \textit{nämlich \((\gamma,\lambda).(\gamma,\lambda^{\prime}),\allowbreak (\gamma,\lambda^{\prime\prime})\) usw. Wir nehmen an, dass in einer Funktion von \(\beta\) Unbestimmten} \textit{von derselben Art wie in Artikel 347, nämlich in \(F=\varphi(t,u,v\ldots)\), für die Unbestimmten \(t,\allowbreak u,\allowbreak v\) usw.} \textit{jeweils die Summen \((\gamma,\lambda),\allowbreak (\gamma,\lambda^{\prime}),\allowbreak (\gamma,\lambda^{\prime\prime})\) usw. eingesetzt werden und ihr Wert nach den Regeln des} \textit{Artikels 345, IV, auf} \[\begin{aligned}
A+a(\gamma,1)+a^{\prime}(\gamma,g)&\ldots+a^\zeta(\gamma,g^{\alpha\beta-\alpha})\\
&\ldots+a^\theta(\gamma,g^{\alpha\beta-1})=W
\end{aligned}\] \textit{zurückgeführt wird. Dann behaupte ich: Ist \(F\) eine symmetrische Funktion, so haben die Perioden in \(W\), die in derselben Periode} \textit{mit \(\beta\gamma\) Gliedern enthalten sind, \textit{das heißt} allgemein solche wie \((\gamma,g^\mu)\) und \((\gamma,g^{\alpha\nu+\mu})\), wobei} \textit{\(\nu\) eine beliebige ganze Zahl bezeichnet, dieselben Koeffizienten.}

\textit{Beweis.} Da die Periode \((\beta\gamma,\lambda g^\alpha)\) mit \((\beta\gamma,\lambda)\) identisch ist, müssen die kleineren Perioden \((\gamma,\lambda g^\alpha),\allowbreak (\gamma,\lambda^{\prime}g^\alpha),\allowbreak (\gamma,\lambda^{\prime\prime}g^\alpha)\) usw., aus denen die erstere offenbar besteht, mit denen übereinstimmen, aus denen die letztere besteht, wenn auch in anderer Reihenfolge. Nimmt man also an, dass nach ihrem jeweiligen Einsetzen für \(t,\allowbreak u,\allowbreak v\) usw. \(F\) in \(W^{\prime}\) übergeht, so stimmt \(W^{\prime}\) mit \(W\) überein. Nach Artikel 347 gilt aber \[\begin{aligned}
W^{\prime}&=A+a(\gamma,g^\alpha)+a^{\prime}(\gamma,g^{\alpha+1})\\
&\qquad\ldots+a^\zeta(\gamma,g^{\alpha\beta})\ldots+a^\theta(\gamma,g^{\alpha\beta+\alpha-1})\\
&=A+a(\gamma,g^\alpha)+a^{\prime}(\gamma,g^{\alpha+1})\\
&\qquad\ldots+a^\zeta(\gamma,1)\ldots+a^\theta(\gamma,g^{\alpha-1})
\end{aligned}\] Da dieser Ausdruck mit \(W\) übereinstimmen muss, muss der erste, zweite, dritte usw. Koeffizient in \(W\) (von \(a\) an gerechnet) jeweils mit dem Koeffizienten an der entsprechenden Stelle aus \(\alpha+1,\allowbreak \alpha+2,\allowbreak \alpha+3\) usw. übereinstimmen. Daraus folgt ohne Mühe, dass allgemein die Koeffizienten der Perioden \((\gamma,g^\mu),\allowbreak (\gamma,g^{\alpha+\mu}),\allowbreak (\gamma,g^{2\alpha+\mu})\ldots(\gamma,g^{\nu\alpha+\mu})\), die an den Stellen \(\mu+1,\allowbreak \alpha+\mu+1,\allowbreak 2\alpha+\mu+1\ldots\nu\alpha+\mu+1\) stehen, einander gleich sein müssen. \textit{Q. E. D.}

Daraus ist klar, dass sich \(W\) auf die Form \[A+a(\beta\gamma,1)+a^{\prime}(\beta\gamma,g)\ldots+a^\varepsilon(\beta\gamma,g^{\alpha-1})\] bringen lässt, wobei alle Koeffizienten \(A,\allowbreak a\) usw. ganze Zahlen sind, wenn alle festen Koeffizienten in \(F\) ganze Zahlen sind. Ferner sieht man leicht: Setzt man danach für die Unbestimmten in \(F\) die \(\beta\) Perioden mit \(\gamma\) Gliedern ein, die in einer anderen Periode mit \(\beta\gamma\) Gliedern, nämlich in \((\beta\gamma,\lambda k)\), enthalten sind und offenbar \((\gamma,\lambda k),\allowbreak (\gamma,\lambda^{\prime}k),\allowbreak (\gamma,\lambda^{\prime\prime}k)\) usw. sind, so ist der daraus hervorgehende Wert \(A+a(\beta\gamma,k)+a^{\prime}(\beta\gamma,gk)\ldots+a^\varepsilon(\beta\gamma,g^{\alpha-1}k)\).
\SourceLeaf{435}

Übrigens ist klar, dass sich der Satz auch auf den Fall erstreckt, in dem \(\alpha=1\) oder \(\beta\gamma=n-1\) gilt; hier sind nämlich \textit{alle} Koeffizienten in \(W\) gleich, sodass sich \(W\) auf die Form \(A+a(\beta\gamma,1)\) zurückführt.


